\paragraph{谐矩—多极层级：双通道粗粒化的主导不变量与高阶不可见性}
定理 \ref{thm:xi-cauchy-poisson-two-channel-orthogonal-leading} 将双通道位置–深度混合的 \(t^{-4}\) 主导项分解为两个正交信息源。
下述结果把该二阶主导模态提升为一条全阶“谐矩—多极”层级：若二阶谐矩被对称性湮灭，则相对熵的主导阶自动上移至更高多极，并且主导常数具有完全闭式。

\begin{theorem}[Poisson--Cauchy 粗粒化相对熵的谐矩—多极主导律]\label{thm:xi-poisson-cauchy-harmonic-moment-multipole-kl}
沿用定理 \ref{thm:xi-cauchy-poisson-two-channel-orthogonal-leading} 的设定，并记
\(
T:=t+\bar\delta
\)。
定义复中心变量与其幂矩
\[
Z:=(\Gamma-\bar\gamma)+i(\Delta-\bar\delta),
\qquad
m_k:=\mathbb E[Z^k]\qquad(k\ge 0),
\]
则 \(m_0=1\)、\(m_1=0\)。
令
\[
k_0:=\min\{k\ge 2:\ m_k\neq 0\},
\]
并假设 \(\mathbb E|Z|^{k_0+1}<\infty\)。
则当 \(t\to\infty\)（等价于 \(T\to\infty\)）时有主导渐近
\begin{equation}\label{eq:xi-poisson-cauchy-harmonic-moment-multipole-kl}
\boxed{\;
D_{\mathrm{KL}}(h_t\mid g_t)
=c_{k_0}\,\frac{|m_{k_0}|^2}{T^{2k_0}}
o\!\left(T^{-2k_0}\right),\;}
\qquad
c_k:=\frac{\binom{2k-2}{k-1}}{2^{2k}}\ (k\ge 2).
\end{equation}
因此，双通道 Poisson--Cauchy 粗粒化相对熵在无穷远尺度的主导阶与主导常数完全由最小非零复中心矩 \(m_{k_0}\) 所锁定；
其系数为中心二项式系数的闭式。
\end{theorem}
\begin{proof}[证明要点（复 resolvent 展开与常数积分）]
令 \(T=t+\bar\delta\) 并记
\(
w:=(x-\bar\gamma)-iT
\)。
由恒等式
\[
\frac{1}{\pi}\frac{s}{x^2+s^2}
=\frac{1}{\pi}\Im\frac{1}{x-is}\qquad(s>0),
\]
在定理 \ref{thm:xi-cauchy-poisson-two-channel-orthogonal-leading} 的定义下有
\[
g_t(x)=\frac{1}{\pi}\Im\frac{1}{w},
\qquad
h_t(x)=\frac{1}{\pi}\Im\,\mathbb E\!\left[\frac{1}{w-Z}\right].
\]
在 \(|Z|\le |w|/2\) 上用几何级数
\(
(1-\frac Zw)^{-1}=\sum_{k\ge 0}(\frac Zw)^k
\)
展开并取期望，结合 \(m_1=0\) 与尾部截断估计（由 \(\mathbb E|Z|^{k_0+1}<\infty\) 控制）得到
\[
h_t(x)=g_t(x)+\frac{1}{\pi}\Im\frac{m_{k_0}}{w^{k_0+1}}+o(T^{-(k_0+1)})
\]
（在足以支撑相对熵二阶展开的意义下）。
令 \(y=(x-\bar\gamma)/T\)，则 \(w=T(y-i)\) 且
\(
g_t(x)=\frac{1}{\pi T}\frac{1}{1+y^2}
\)。
记
\(
\delta_t:=\frac{h_t-g_t}{g_t}
\)，
则
\[
\delta_t(x)
=T^{-k_0}\Im\!\left(m_{k_0}\frac{1+y^2}{(y-i)^{k_0+1}}\right)
o(T^{-k_0})
=T^{-k_0}\Im\!\left(m_{k_0}\frac{y+i}{(y-i)^{k_0}}\right)
o(T^{-k_0}).
\]
由于 \(\int g_t\delta_t=\int(h_t-g_t)=0\) 且 \(\delta_t\to 0\)，
由二阶展开
\(
(1+u)\log(1+u)=u+\frac12u^2+o(u^2)
\)
得到
\[
D_{\mathrm{KL}}(h_t\mid g_t)=\frac12\int_{\RR} g_t\,\delta_t^2\,\dd x+o(T^{-2k_0}).
\]
代入 \(x=\bar\gamma+Ty\) 并用奇偶性消去交叉项，归结为
\[
D_{\mathrm{KL}}(h_t\mid g_t)
=\frac{|m_{k_0}|^2}{2T^{2k_0}}
\int_{\RR}\frac{1}{\pi}\frac{1}{1+y^2}\Bigl(\Im\frac{y+i}{(y-i)^{k_0}}\Bigr)^2\dd y
o(T^{-2k_0}).
\]
最后计算积分常数。令 \(y=\tan\theta\)（\(\theta\in(-\pi/2,\pi/2)\)），则 \(\dd y/(1+y^2)=\dd\theta\) 且可化简为
\[
\frac{y+i}{(y-i)^{k}}
=\cos^{k-1}\theta\cdot e^{i(k+1)(\frac{\pi}{2}-\theta)}.
\]
故
\(
\bigl(\Im\frac{y+i}{(y-i)^{k}}\bigr)^2
=\cos^{2k-2}\theta\cdot \sin^2((k+1)(\frac{\pi}{2}-\theta))
\)。
将 \(\sin^2=\frac12(1-\cos(2\cdot))\) 展开并用 \(\cos^{2k-2}\theta\) 的 Fourier 频率界与正交性得到
\[
\int_{-\pi/2}^{\pi/2}\cos^{2k-2}\theta\sin^2((k+1)(\tfrac{\pi}{2}-\theta))\,\dd\theta
=\frac12\int_{-\pi/2}^{\pi/2}\cos^{2k-2}\theta\,\dd\theta
=\frac{\pi}{2}\frac{\binom{2k-2}{k-1}}{2^{2k-2}}.
\]
代回即得
\(
\int_{\RR}\frac{1}{\pi}\frac{1}{1+y^2}\bigl(\Im\frac{y+i}{(y-i)^{k}}\bigr)^2\dd y
=\binom{2k-2}{k-1}/2^{2k-1}
\)，
从而 \eqref{eq:xi-poisson-cauchy-harmonic-moment-multipole-kl} 成立。
可复现核验脚本：\texttt{scripts/exp\_xi\_poisson\_cauchy\_harmonic\_multipole\_kl\_audit.py}。证毕。
\end{proof}

\begin{remark}[二阶谐矩的协方差谱隙表达]\label{rem:xi-poisson-cauchy-m2-covariance-anisotropy}
在定理 \ref{thm:xi-poisson-cauchy-harmonic-moment-multipole-kl} 的记号下，
二阶复中心矩可写为
\[
m_2=\mathbb E[Z^2]=(\sigma_\gamma^2-\sigma_\delta^2)+2i\,\sigma_{\gamma\delta},
\]
从而
\(
|m_2|^2=(\sigma_\gamma^2-\sigma_\delta^2)^2+4\sigma_{\gamma\delta}^2
\)。
因此当 \(k_0=2\) 时，\eqref{eq:xi-poisson-cauchy-harmonic-moment-multipole-kl} 与
\eqref{eq:xi-cauchy-poisson-two-channel-kl-leading} 的 \(T^{-4}\) 主导项完全一致。
\end{remark}

\begin{proposition}[一般 \texorpdfstring{$f$}{f}-散度的同型多极主导律]\label{prop:xi-poisson-cauchy-fdiv-harmonic-multipole}
在定理 \ref{thm:xi-poisson-cauchy-harmonic-moment-multipole-kl} 的条件下，
令 \(f:(0,\infty)\to\RR\) 为凸函数，满足 \(f(1)=0\)，并在 \(1\) 的邻域二阶可微且 \(f''(1)<\infty\)。
定义 Csisz\'ar \texorpdfstring{$f$}{f}-散度
\[
D_f(h\mid g):=\int_{\RR} g(x)\,f\!\left(\frac{h(x)}{g(x)}\right)\,\dd x.
\]
则当 \(t\to\infty\) 时有
\[
\boxed{\;
D_f(h_t\mid g_t)
=f''(1)\,c_{k_0}\,\frac{|m_{k_0}|^2}{T^{2k_0}}
o\!\left(T^{-2k_0}\right)\;}
\qquad(T=t+\bar\delta),
\]
其中常数 \(c_k\) 同 \eqref{eq:xi-poisson-cauchy-harmonic-moment-multipole-kl}。
\end{proposition}
\begin{proof}[证明要点]
令 \(r_t=h_t/g_t=1+\delta_t\)。
由定理 \ref{thm:xi-poisson-cauchy-harmonic-moment-multipole-kl} 的推导，\(\delta_t\to0\) 且 \(\int g_t\delta_t=0\)。
对 \(f(1+u)\) 在 \(u=0\) 处作二阶 Taylor 展开并积分，线性项消失，得到
\(
D_f(h_t\mid g_t)=\frac{f''(1)}{2}\int g_t\delta_t^2+o(\int g_t\delta_t^2)
\)。
而 \(\int g_t\delta_t^2\) 的主项由同一定理中的积分常数给出，从而得到结论。证毕。
\end{proof}

\begin{corollary}[相对熵耗散量的多极幂律常数]\label{cor:xi-poisson-cauchy-fisher-harmonic-multipole}
在定理 \ref{thm:xi-poisson-cauchy-harmonic-moment-multipole-kl} 的条件下，
由命题 \ref{prop:xi-poisson-kl-dissipation-fractional-fisher} 的耗散恒等式得
\[
\boxed{\;
\mathcal I_1(h_t\mid g_t)
=2k_0\,c_{k_0}\,\frac{|m_{k_0}|^2}{T^{2k_0+1}}
o\!\left(T^{-(2k_0+1)}\right)\;}
\qquad(T=t+\bar\delta).
\]
\end{corollary}
\begin{proof}
对 \eqref{eq:xi-poisson-cauchy-harmonic-moment-multipole-kl} 关于 \(t\) 求导并用
\(
\frac{\dd}{\dd t}D_{\mathrm{KL}}(h_t\mid g_t)=-\mathcal I_1(h_t\mid g_t)
\)
即可。证毕。
\end{proof}

\begin{proposition}[任意阶多极不可见性的显式构造与主部常数]\label{prop:xi-poisson-cauchy-multipole-invisibility-construction}
对任意整数 \(N\ge 2\) 与任意 \(r>0\)，可取一个有限支撑概率测度 \(\tilde\nu\)（支撑于 \(\RR\times(0,\infty)\)）
使其复中心矩满足
\[
m_2=\cdots=m_{N-1}=0,
\qquad
m_N=r^N\neq 0.
\]
对该 \(\tilde\nu\) 诱导的 \(h_t,g_t\)（定理 \ref{thm:xi-cauchy-poisson-two-channel-orthogonal-leading}），有
\[
D_{\mathrm{KL}}(h_t\mid g_t)
=c_N\,\frac{r^{2N}}{T^{2N}}+o(T^{-2N}),
\qquad
c_N=\frac{\binom{2N-2}{N-1}}{2^{2N}}.
\]
特别地取 \(N=4\) 得
\(
D_{\mathrm{KL}}(h_t\mid g_t)=\frac{5}{64}\,r^8\,T^{-8}+o(T^{-8})
\)。
\end{proposition}
\begin{proof}[证明要点]
取 \(\omega_j=e^{2\pi ij/N}\) 并令 \(Z_j=r\omega_j\)（\(0\le j\le N-1\)）。
取任意 \(\bar\delta>r\)，并定义有限支撑测度
\[
\tilde\nu:=\frac{1}{N}\sum_{j=0}^{N-1}\delta_{(\bar\gamma+\Re Z_j,\ \bar\delta+\Im Z_j)}.
\]
则 \(\Delta-\bar\delta\in[-r,r]\) 且 \(\Delta>0\)。
并且对 \(1\le k\le N-1\) 有
\(
m_k=\frac{r^k}{N}\sum_{j=0}^{N-1}\omega_j^k=0
\)，而 \(m_N=r^N\)。
于是 \(k_0=N\)，代入定理 \ref{thm:xi-poisson-cauchy-harmonic-moment-multipole-kl} 即得结论。证毕。
\end{proof}

\endinput

